problem:  
Let \((X,d,\mathfrak{m})\) be a compact metric measure space satisfying the curvature–dimension condition \(CD(K,N)\) for some \(K\in\mathbb{R}\), \(1<N<\infty\).  
Assume:  
1. For each \(x\in X\), every measured tangent cone \(T_xX\) (in the pointed measured Gromov–Hausdorff sense) is uniquely Euclidean of the same dimension \(k_0\).  
2. The canonical heat flow \((P_t)_{t>0}\) on \(L^2(X,\mathfrak{m})\), regarded as the \(L^2\)-gradient flow of the Cheeger energy \(\mathrm{Ch}\), admits a transition kernel \(p_t(x,y)\) with short-time Gaussian two-sided bounds and asymptotically Euclidean form  
   \[
   p_t(x,y)\approx (4\pi t)^{-k_0/2}e^{-d(x,y)^2/(4t)}(1+o(1))+ \text{terms uniform in small balls.}
   \]
3. No infinitesimal Hilbertianity (quadraticity of \(\mathrm{Ch}\)) is assumed a priori.  

**Question:**  
Do these assumptions already force the Cheeger energy \(\mathrm{Ch}\) to be quadratic—thereby implying that \((X,d,\mathfrak{m})\) is an \(RCD(K,N)\) space?  
Equivalently, can precise Gaussian short-time heat kernel asymptotics combined with globally Euclidean tangent cones alone enforce genuinely Riemannian (Hilbertian) first-order calculus, ruling out \(CD(K,N)\) spaces of Finsler type?  

Why is it a "good" problem:  
This question eliminates known equivalences such as the Bakry–Émery gradient estimate and focuses squarely on the **heat kernel as an analytic fingerprint** of Riemannian structure. It asks whether the precise form of small-time diffusion behavior already encodes quadratic energy, thus bridging three core aspects of modern geometric analysis: synthetic curvature dimension theory, Dirichlet form and diffusion analysis, and tangent cone geometry.  

A positive resolution would offer a striking analytic characterization of the \(RCD\) property—showing that sharp Euclidean short-time diffusion asymptotics suffice to recover infinitesimal Hilbertianity. A counterexample would reveal a new class of exotic “Finsler-type” \(CD(K,N)\) spaces with deceptively Gaussian heat propagation. The problem is simple yet profound: can one *hear* the Riemannian metric from the sound of diffusion? Its resolution would unite geometric measure structure and PDE asymptotics in the nonsmooth setting, making it an exemplary "narrow entrance, profound depth" research problem.